\section{任意の重みで使える標準部分の還元}
\label{sec:notes-all-weight-filtered-v40}

\paragraph{固定する対象。}
$m=2j+r\ge3$, $j,r\ge0$、$\alpha=0,s=1$ のCGG座標代数の
$(0,1,m,m+1)$ 切出しを $B_{j,r}$、純粋三変数部分を $N_m$ とする。
$r=0$ を含む。$n_d=\binom{d+2}{2}$ と書く。
任意の標識付きArtin変形で $N_m$ は自由直和因子の部分代数として一意に持ち上がり、
二段・三段合成の像の階数は $n_m,n_m,n_{m+1}$ に保たれる。
同型と、冪等元を零にする正規化導分もこの部分を保つ。

\paragraph{行列計算から独立した消滅機構。}
可換積の正の $w$ 重みの二cocycleに対し、中間の対称積の $m-3$ 因子を固定する。
残る四変数への二重反対称化は重み三と同じで、係数 $n(n-2)$ により一次・二次が消滅する。
CGG積では補助重み $(0,1,2,m+r+1)$ を使う。$N$ は一つ下げ、$X_0$ は保つので、
有限フィルトレーションのassociated gradedは可換の複体である。
従って $Q=C/C_F$ に対し
\[
 H^1(Q)=H^2(Q)=0,\qquad \dim H^3(Q)=4m^2-6m-7.
\]
この二段階の証明は、非零量子点に個別の行基本変形を繰り返す必要をなくす。

\paragraph{相対安定化群を殺す明示的な汎関数。}
純粋部分を固定する正規化導分は、二つの原始射で共通に $w\mapsto aw+p$ と書ける。
$q,y$ との混合関係により
$Np=aq^jz^r$、$X_0p=(m+r+1)p$ である。
許される $p$ の基底は
$x^b y^{2j-1-2b}z^{b+r+1}$、$0\le b<j$。
$N$ の二重対角行列を左から零化する汎関数は
\[
 \ell_0=1,\qquad
 \ell_{b+1}=-\frac{2(b+r+1)}{2j-1-2b}\ell_b.
\]
$q^jz^r$ とのpairingは
$(-1)^j\sum_{b=0}^j|\ell_b|\binom jb\ne0$ だから $a=0$、続いて $p=0$ となる。
これは全共鳴に共通する有限次元の障害を、具体的な有理汎関数で表示している。

\paragraph{相対変形の次元と通常hull。}
純粋部分の変形とその同定を固定した相対変形は滑らかで
\[
 k_m=3n_{m+1}-2n_m+7=\frac{m^2+9m+28}{2}
\]
次元である。原始射を含む三重合成がないため相対結合性条件はなく、
上記の汎関数が相対自己同型を消す。
最初の値は $k_3=32,k_4=40,k_5=49,k_6=59,k_7=70$。
通常の変形接空間については
\[
 0\to HH^1(B_{j,r})\to HH^1(N_m)
 \xrightarrow{\partial}\mathbf C^{k_m}
 \to HH^2(B_{j,r})\to HH^2(N_m)\to0
\]
を使う。今回、$\partial=0$ も全ての重みで証明できた。
従って適合するminiversal環について
\[
 R_{B_{j,r}}\simeq R_{N_m}[[u_1,\ldots,u_{k_m}]]
\]
となる。$N_m$ と $k_m$ は $j,r$ によらないので、固定した重みの全共鳴は
端点も含め同じ通常変形環を持つ。対称性保持部分まで同じになるという主張ではない。

\paragraph{純粋部分の外的導分を一様に計算する方法。}
可換点の導分で三段積 $vtz=ztv$ を微分すると、左右の次数一作用の差はスカラーになる。
自然な $\mathfrak{gl}_3$ 作用と内的導分を引いた後、中間作用は
$xG(yt)=yG(xt)$ を満たすのでスカラーである。よって可換点の $HH^1$ は $\mathfrak{pgl}_3$。
三本の $x,x^{m-1},x$ を可逆にする普遍局所化は $M_4(\mathbf C[Y,Z])$ である。
$m\ge3$ が、$yzx^{m-3}$ から $YZ=ZY$ を導くために必要となる。
CGG族を $\mathbf C[h]/h^n$ 上で局所化すると、既知の平坦な座標chartへの全射は法 $h$ で同型。
核 $K$ は、平坦な商へのTor完全列から $K/hK=0$、従って $K=0$ になる。
元の局所化の平坦性を先に仮定しないことが重要である。
一次で持ち上がる導分は $2(Z-Y^2)\partial_Y\wedge\partial_Z$ を保つ。
射影的ベクトル場の中でその安定化群は $\partial_Y+2Y\partial_Z$ の一次元である。
頂点内的導分の像は三本の隣接射の自己係数により自由直和因子と分かる。
その商での飽和された導分核への特殊化と、実際の $h^{-1}\log(1+hN)$ により
$HH^1(N_m)=\mathbf C$ が全ての重みで従う。
その生成元は四変数側にも存在するため、上の接続写像 $\partial$ は零である。

\paragraph{次の比較問題。}
純粋部分の二次Hochschild群と全変形環を非可換射影平面のものと比較することが有望である。
独立な有理数計算では $m=3,4,5$ の $N_m$ について
\[
\begin{array}{c|c|c|c}
 m & (HH^1,HH^2)\text{、可換点}&(HH^1,HH^2)\text{、非零CGG点}&\dim HH^3\\\hline
 3&(8,10)&(1,3)&28\\
 4&(8,10)&(1,3)&108\\
 5&(8,10)&(1,3)&270
\end{array}
\]
となる。一次群の全重みの主張は上で証明済みだが、表の二次群と滑らかさを一般の $m$ へ外挿しない。
特に $(\mathcal O,\mathcal O(1),\mathcal O(m),\mathcal O(m+1))$ は
$\mathbb P^2$ 上では逆向きの高次Extを持つので、通常のtilting同値をそのまま使えない。
一般比較を証明できれば、未決の $R_{N_m}$ 自体と全cornerの通常hullを決定できる。
対称性保持曲線が原始射の接合によって何次に切り詰められるかは、その後にも残る別の問題である。

\paragraph{検証の範囲。}
通常ゲートは $3\le m\le7$ の全十六共鳴と、五つの可換正重み複体、
七十七の極端重量の汎関数、および射影的Poisson安定化群の階数七を厳密に確認する。
全整数の証明は上記の構造による。純粋部分の表は別の任意実行プローブで再現でき、
一般定理の証拠へ読み替えない。
